% -----------------------------------------------------------------------------
% -----------------------------------------------------------------------------

\begin{exbox}[Compton Ward 恒等式]
Compton 振幅~\eqnrefs{eq:compton-M} 含两张树图。规范不变性的直接检验不是“每张图都横向”，而是把某一条外光子的偏振矢量替换为它自己的四动量，检查\emph{两图之和}是否为零。以入射光子为例，令
\begin{equation*}
 \slashed\varepsilon\longrightarrow\slashed k.
\end{equation*}
第一张图的电子传播子动量为 $r_s=p+k$。利用
\begin{equation*}
 \slashed k
 =(\slashed r_s-\me c)-(\slashed p-\me c),
 \qquad
 (\slashed p-\me c)u(p)=0,
\end{equation*}
可把分子与逆传播子配成
\begin{align*}
 (\slashed r_s+\me c)\slashed k\,u(p)
 &=(\slashed r_s+\me c)(\slashed r_s-\me c)u(p),\\
 &=(r_s^2-\me^2c^2)u(p),\\
 &=2p\cdot k\,u(p).
\end{align*}
所以第一张图在 Ward 替换后只剩
\begin{equation*}
 \bar u(p')\slashed\varepsilon'^{*}u(p).
\end{equation*}

第二张图令 $r_u=p-k'=p'-k$。从左侧使用出射电子 Dirac 方程，
\begin{align*}
 \bar u(p')\slashed k(\slashed r_u+\me c)
 &=\bar u(p')(\slashed p'-\slashed r_u)(\slashed r_u+\me c),\\
 &=\bar u(p')(\me c-\slashed r_u)(\slashed r_u+\me c),\\
 &=\bar u(p')(\me^2c^2-r_u^2),\\
 &=2p\cdot k'\,\bar u(p').
\end{align*}
因此第二张图在\eqnrefs{eq:compton-M} 自带的相对负号作用下给
\begin{equation*}
 -\bar u(p')\slashed\varepsilon'^{*}u(p).
\end{equation*}
两项精确相消：
\begin{equation}
 \boxed{
 \mathcal M_C(\varepsilon\to k)=0.}
 \label{eq:compton-ward-incoming-dp64}
\end{equation}
对出射光子作 $\varepsilon'^{*\mu}\to k'^{\mu}$ 的同样计算也得到零。因此若把振幅写成
\begin{equation*}
 \mathcal M_C
 =\varepsilon'^{*}_{\mu}\varepsilon_\nu\,
 \mathcal T^{\mu\nu},
\end{equation*}
就有
\begin{equation}
 \boxed{
 k_\nu\mathcal T^{\mu\nu}=0,
 \qquad
 k'_\mu\mathcal T^{\mu\nu}=0.}
 \label{eq:compton-double-ward-dp64}
\end{equation}

这也把“规范变换不影响散射振幅”变成一个一行可检验的结论。因为真实光子的规范代表元可改变为
\begin{equation*}
 \varepsilon^\mu\to\varepsilon^\mu+\lambda k^\mu,
\end{equation*}
振幅变化量正比于 $k_\nu\mathcal T^{\mu\nu}$，所以为零。辅助参考矢量 $n^\mu$ 出现在物理偏振完备张量中时，所有含 $k^\mu$ 或 $k'^\mu$ 的项也会被\eqnrefs{eq:compton-double-ward-dp64} 消掉。于是未分辨光子偏振时，可以在\emph{与完整守恒振幅收缩以后}安全使用
\begin{equation*}
 \sum_\lambda\varepsilon_\lambda^\mu\varepsilon_\lambda^{\nu *}
 \rightsquigarrow-\eta^{\mu\nu}.
\end{equation*}
偏振求和必须作用在满足 Ward 条件的完整树级振幅上；单独一张 Compton 图不满足该横向性条件。


\medskip

\eqnrefs{eq:compton-M} 是同一电子线上两个光子顶角的两种时间次序之和。单独一张图并不满足外光子的横向性；需要把偏振矢量替换成相应光子四动量，并利用外电子的在壳 Dirac 方程，才能看到两张图的非横向部分恰好相消。先得到入射光子的 Ward 条件~\eqnrefs{eq:compton-ward-incoming-dp64}，再对出射光子重复同一代数，最终得到双横向条件~\eqnrefs{eq:compton-double-ward-dp64}。


把\eqnrefs{eq:compton-M} 写成
\begin{equation*}
 \mathcal M_C=-g_{\rm em}^2\bar u(p')
 \left[A_s-A_u\right]u(p),
\end{equation*}
其中
\begin{align*}
 A_s&=\slashed\varepsilon'^{*}
 \frac{\slashed p+\slashed k+\me c}{2p\cdot k}
 \slashed\varepsilon,\\
 A_u&=\slashed\varepsilon
 \frac{\slashed p-\slashed k'+\me c}{2p\cdot k'}
 \slashed\varepsilon'^{*}.
\end{align*}
先检验入射 Ward 条件。令 $r_s=p+k$，则
\begin{equation*}
 r_s^2-\me^2c^2
 =p^2+2p\cdot k+k^2-\me^2c^2
 =2p\cdot k.
\end{equation*}
在 $\slashed\varepsilon\to\slashed k$ 后，
\begin{align*}
 A_su(p)
 &=\slashed\varepsilon'^{*}
 \frac{\slashed r_s+\me c}{2p\cdot k}
 \left[(\slashed r_s-\me c)-(\slashed p-\me c)\right]u(p),\\
 &=\slashed\varepsilon'^{*}
 \frac{(\slashed r_s+\me c)(\slashed r_s-\me c)}{2p\cdot k}u(p),\\
 &=\slashed\varepsilon'^{*}u(p).
\end{align*}
再令 $r_u=p-k'=p'-k$。由 $\bar u(p')(\slashed p'-\me c)=0$，
\begin{align*}
 \bar u(p')\slashed k(\slashed r_u+\me c)
 &=\bar u(p')(\slashed p'-\slashed r_u)(\slashed r_u+\me c),\\
 &=\bar u(p')(\me c-\slashed r_u)(\slashed r_u+\me c),\\
 &=\bar u(p')(\me^2c^2-r_u^2).
\end{align*}
另一方面
\begin{equation*}
 r_u^2-\me^2c^2
 =(p-k')^2-\me^2c^2
 =-2p\cdot k',
\end{equation*}
所以 $\me^2c^2-r_u^2=2p\cdot k'$，从而
\begin{equation*}
 \bar u(p')A_u u(p)
 \xrightarrow{\varepsilon\to k}
 \bar u(p')\slashed\varepsilon'^{*}u(p).
\end{equation*}
由于完整振幅中是 $A_s-A_u$，两项相消，得到\eqnrefs{eq:compton-ward-incoming-dp64}。

出射 Ward 条件的代数完全平行。把 $\slashed\varepsilon'^{*}\to\slashed k'$，第一项从左侧利用
\begin{equation*}
 k'=r_s-p',
 \qquad
 \bar u(p')(\slashed p'-\me c)=0,
\end{equation*}
第二项从右侧利用
\begin{equation*}
 k'=p-r_u,
 \qquad
 (\slashed p-\me c)u(p)=0.
\end{equation*}
两张图分别化为同一个 $\bar u(p')\slashed\varepsilon u(p)$，再由相对负号相消，因此
\begin{equation*}
 \mathcal M_C(\varepsilon'^{*}\to k')=0.
\end{equation*}
把两条结果从任意外部偏振因子中剥离，就得到\eqnrefs{eq:compton-double-ward-dp64}。

最后取一般规范代表元
\begin{equation*}
 \varepsilon^\nu\to\varepsilon^\nu+\lambda k^\nu,
 \qquad
 \varepsilon'^{*\mu}\to\varepsilon'^{*\mu}+\lambda' k'^{\mu}.
\end{equation*}
振幅的线性变化为
\begin{equation*}
 \delta\mathcal M_C
 =\lambda\varepsilon'^{*}_{\mu}k_\nu\mathcal T^{\mu\nu}
 +\lambda'k'_\mu\varepsilon_\nu\mathcal T^{\mu\nu}
 +\lambda\lambda'k'_\mu k_\nu\mathcal T^{\mu\nu}=0.
\end{equation*}
所以 Ward 恒等式既解释了规范变换不改变振幅，也解释了为什么物理偏振投影中的参考矢量依赖在可观测量中消失。
\end{exbox}


